% Pseudocódigos e detalhamento de implementação — embutido no relatório

\section{Pseudocódigos complementares dos analisadores}
\label{app:pseudocode}

\subsection{Analisador léxico por varredura com reconhecimento de padrões}

O analisador léxico implementado em \texttt{lexer.py} adota varredura linear sobre o código-fonte com funções auxiliares de reconhecimento. Embora a especificação acadêmica tradicional empregue autômatos finitos, a implementação atual privilegia clareza didática e extensibilidade para tokens orientados a objetos e de paralelismo estruturado.

\begin{lstlisting}[caption={Pseudocodigo do analisador lexico}]
funcao tokenize(fonte):
    posicao <- 0
    enquanto posicao < comprimento(fonte):
        ignorar espacos em branco
        se comentario detectado entao avancar e continuar
        se digito entao emitir NUMBER_LITERAL
        senao se aspas duplas entao emitir STRING_LITERAL
        senao se letra ou sublinhado entao
            lexema <- ler identificador ou palavra-chave
            emitir tipo correspondente em KEYWORDS ou IDENTIFIER
        senao
            emitir operador ou delimitador de tabela fixa
    emitir EOF
    retornar lista de tokens
\end{lstlisting}

\subsection{Analisador sintático descendente recursivo}

O analisador sintático consome a lista de tokens produzida na fase anterior e constrói a árvore sintática abstrata por meio de um conjunto de funções recursivas, uma para cada não-terminal da gramática BNF. A precedência de expressões materializa-se na cadeia de chamadas de \texttt{assignment}, \texttt{logical\_or}, \texttt{additive} e demais níveis, conforme a gramática da Seção de especificação da linguagem.

\begin{lstlisting}[caption={Pseudocodigo simplificado de expressoes com precedencia}]
funcao parse_expr():
    retornar parse_assignment()

funcao parse_assignment():
    no <- parse_conditional()
    se token atual for "=" entao
        consumir "="
        direita <- parse_assignment()
        retornar no de atribuicao
    retornar no

funcao parse_additive():
    esquerda <- parse_multiplicative()
    enquanto token em {+, -}:
        operador <- consumir operador
        direita <- parse_multiplicative()
        esquerda <- no binario(esquerda, operador, direita)
    retornar esquerda
\end{lstlisting}

\subsection{Analisador semântico e tabela de símbolos hierárquica}

O visitador semântico percorre a árvore em profundidade, mantendo pilha de escopos. Cada declaração de classe registra símbolo no escopo global; membros de classe abrem escopo filho. A verificação de herança ocorre em passagem posterior, após o registro de todas as classes, garantindo que a superclasse referenciada em \texttt{extends} já conste da tabela.

\begin{lstlisting}[caption={Pseudocodigo do visitador semantico}]
funcao visitar_programa(no):
    para cada declaracao em no.declarations:
        visitar(declaracao)
    para cada classe com extends:
        se superclasse nao registrada entao
            registrar erro semantico

funcao visitar_classe(no):
    se no.nome ja existe entao registrar erro de duplicata
    inserir simbolo de classe no escopo global
    abrir escopo de classe
    para cada membro em no.members:
        se nome de membro repetido entao registrar erro
        visitar(membro)
    fechar escopo de classe
\end{lstlisting}

\subsection{Geração de código de três endereços}

O gerador de representação intermediária percorre a árvore e emite instruções com temporários monotonicamente numerados. Chamadas de método em contexto orientado a objetos recebem prefixo com o nome da classe para evitar colisão de símbolos na fase de emissão para linguagem C.

\begin{lstlisting}[caption={Pseudocodigo de lowering para TAC}]
funcao gerar_programa(no):
    para cada declaracao em no.declarations:
        gerar(declaracao)
    emitir CALL_MAIN

funcao gerar_method(no, classe_atual):
    nome_funcao <- classe_atual + "_" + no.nome
    emitir FUNC_BEGIN, nome_funcao
    para cada parametro em no.parameters:
        emitir FUNC_PARAM, parametro.nome
    gerar(no.corpo)
    emitir FUNC_END, nome_funcao

funcao gerar_par(no):
    emitir PAR_BEGIN
    para cada comando em no.statements:
        gerar(comando)
    emitir PAR_END
\end{lstlisting}

\subsection{Emissão de código C e invocação do compilador externo}

O back-end em linguagem C percorre a lista de instruções de três endereços e produz um arquivo de texto com funções auxiliares e ponto de entrada \texttt{main}. O método \texttt{finalize} grava o artefato em arquivo temporário, invoca \texttt{gcc} com otimização de segundo nível e captura a saída padrão e o código de retorno do processo filho.

\begin{lstlisting}[caption={Pseudocodigo do back-end C}]
funcao emit(ast_serializada):
    tac <- TACGenerator.lower(ast_serializada)
    codigo_c <- SimpleCCodeGenerator.generate(tac)
    armazenar codigo_c internamente

funcao finalize():
    arquivo <- gravar temporario(codigo_c)
    resultado_compilacao <- executar(["gcc", "-O2", "-lm", arquivo])
    se codigo de retorno diferente de zero entao
        retornar TranslationResult com mensagem de erro
    saida <- executar(binario gerado)
    retornar TranslationResult(saida, codigo_c, codigo de retorno)
\end{lstlisting}

\section{Registro formal dos \emph{sprints}}
\label{app:sprints}

A Tabela~\ref{tab:sprint-backlog} consolida as entregas planejadas e realizadas em cada iteração, com responsabilidade compartilhada pela equipe do projeto integrado.

\begin{table}[htbp]
\centering
\small
\caption{Registro de entregas por sprint}
\label{tab:sprint-backlog}
\begin{tabularx}{\textwidth}{|c|>{\raggedright\arraybackslash}p{2.2cm}|X|>{\raggedright\arraybackslash}p{2.8cm}|}
\hline
\textbf{Sprint} & \textbf{Período} & \textbf{Entregas principais} & \textbf{Responsáveis} \\
\hline
1 & 2 a 4 de junho de 2026 & Núcleo \texttt{minipar-core}; microsserviços de análise; AST JSON; erros sintático e semântico & Equipe Compiladores \\
2 & 3 a 5 de junho de 2026 & Módulo \texttt{translation}; interpretador; geradores C, Rust e ARM; orquestração HTTP & Equipe Reuso \\
3 & 5 a 7 de junho de 2026 & Coordenador distribuído; workers socket; fractal Sierpinski; extensão Python & Equipe LPS \\
\hline
\end{tabularx}
\end{table}
